\subsection{Controlled Grounding Diagnostics and Numerical Correctness}
\label{subsec:grounding-followup}

Coarse readiness and whole-query numerical correctness rank the existing
pipelines differently. Saved records yield 40/331 full-gold-key scalar
tolerance successes for constrained assignment versus 7/331 for typed greedy,
with 36 gains and three losses. The other matched all-key counts are 37/331 for maximum-weight matching, 35/331 for role repair, 46/331 for search-structured and 9/331 for semantic-IR repair (Table~\ref{tab:nlp4lp-grounding-baselines}). The constrained eligible-only count is 41:
query 102 covers only three of seven gold scalar keys and is excluded from
the all-key count. Candidate construction, context, scoring, abstention and
decoding differ together; both methods already prohibit mention reuse.
Consequently, the gain cannot be attributed solely to global search.

A separate retrospective factorial study uses 49 primary queries and 277
gold slots, holding each candidate pool and score matrix fixed within
decoder comparisons. The 64 collected cases additionally contain 12
development cases, one historical-overlap reconciliation case, and two
reference-only cases. Inference uses predicted-schema scalar metadata
without gold-key filtering, but still receives reference-derived slot structure. Under contextual scores and one-to-one
assignment, exact slot counts increase from 106 to 136 for typed candidates
and from 100 to 133 for rich candidates, but whole-query tolerance successes
remain 5/49. These decoder contrasts do not survive the separately retained
Holm family of 240 exploratory endpoint tests. With reusable mentions and
additive scores, independent selection is already globally optimal.

The earlier joint-grounding pilot likewise finds 3/49 successes for both
local and joint selection on the overlap-free slice. Of the original 50
cases, 32 have a criterion-correct mapping in the common pools, but only
eight retain one in the beam's top ten and three select it. Exact diagnosis
finds 23 strict objective-preference failures and six tied-score ambiguities
among those 32. A subsequent exact value/type-distinct enumeration on the
same 49 primary cases leaves retention at 9/49 and selected correctness at
3/49. More search or numerical diversity therefore produced no observed
whole-instance improvement for these frozen objectives. These are conditional negative findings on a small pilot, not an equivalence result or evidence that better-supported constraints cannot help.

Appendix~\ref{S-app:grounding-diagnostics} retains the full factorial and
pruning evidence, reference-convention corrections and reproduction map.
All these comparisons are retrospective; the known query-291 overlap is
excluded from new primary comparisons, but prior benchmark exposure remains.
